\documentclass[12pt]{article}
\usepackage[utf8]{inputenc}
\usepackage[T1]{fontenc}
\usepackage{hyperref}
\title{Kamakan Language Resources for K-12 Curriculum Development}
\author{Compiled by Claude Sonnet 5 (claude-sonnet-5) for Andr\'as Kornai}
\date{}
\begin{document}
\maketitle
\section{Introduction}
Kamakan (ISO 639-3 vkm; Kamakanan family, Macro-Je, Bahia, Brazil) is extinct. Sources disagree on the date, from the first half of the twentieth century to the 1950s or 1970. The descendants are reported to have merged into the Pataxo-ha-ha-hae population. No native-language curriculum materials were found, and a curriculum would have to be a heritage and revitalization programme built on very thin historical documentation.
\section{Primary Educational Resources}
None found in search results. No primers, textbooks or school programmes for Kamakan were found.
\section{Linguistic and Reference Materials}
\begin{itemize}
\item Wikipedia overview of the Kamak\~a language: \url{https://en.wikipedia.org/wiki/Kamak%C3%A3_language}
\item Wikipedia overview of the Kamak\~a language family: \url{https://en.wikipedia.org/wiki/Kamak%C3%A3_languages}
\item Meni\'en language (related): \url{https://en.wikipedia.org/wiki/Meni%C3%A9n_language}
\item Masakar\'a language (related): \url{https://en.wikipedia.org/wiki/Masakar%C3%A1_language}
\item Kotox\'o language (related): \url{https://en.wikipedia.org/wiki/Kotox%C3%B3_language}
\item John Carter Brown Library Indian Languages Database, which lists two historical Kamakan vocabularies (Meniengs, 90 terms; Mongoyos, 150 terms, in French translation by Eyri\`es): \url{https://www.brown.edu/Facilities/John_Carter_Brown_Library/exhibitions/ildb/details.php?id=555}
\item ASJP wordlist for Kamakan: \url{https://asjp.clld.org/languages/KAMAKAN}
\end{itemize}
\section{Religious and Community Resources}
None found.
\section{Digital and Audio Resources}
No audio or software was found. The historical wordlists above are the only primary data seen.
\section{Curriculum Development Context}
With no speakers and only short historical wordlists, any curriculum would be heritage-oriented, covering history, the Bahia peoples and the little language data that survives. Any use of the language itself should be developed with the descendant community.
\end{document}